\subsection{\problemblock{*Wind} Data Block}\label{sec:block-wind}

Atmospheric wind relative to the earth's surface is defined with a
\problemblock{*wind} block. If no block is supplied, all wind velocities are zero.
Wind is divided into a vertical component and a horizontal component. Horizontal
wind may be specified by speed and heading or by east and north components, and
either the local geocentric or local geodetic horizon system may be selected.

Wind heading is the direction \emph{from which} the wind blows, measured
clockwise from north as shown in \cref{fig:wind-heading-direction}. Thus a heading
of $180^\circ$ means wind from the south, moving north. The east component is
positive for wind from the east, moving west, and the vertical component is positive
for downward air motion toward the ground.

\begin{figure}[htbp]
  \centering
  \begin{tikzpicture}[>=Latex,font=\small,scale=0.95]
  \coordinate (O) at (0,0);
  \draw[->] (O) -- (0,2.7) node[above]{North};
  \draw[->] (O) -- (3.2,0) node[right]{East};
  \draw[->,very thick] (2.5,2.15) -- (0.55,0.45) node[midway,above right]{Wind vector};
  \draw[->] (0,1.0) arc[start angle=90,end angle=40,radius=1.0];
  \node at (0.95,1.15) {Heading};
\end{tikzpicture}

  \caption{Wind horizontal-component direction.}
  \label{fig:wind-heading-direction}
\end{figure}

The block begins with \problemblock{*wind}, followed by \texttt{geocentric} or
\texttt{geodetic}. Three variables and values are then selected from one of the two
consistent sets in \cref{tab:wind-data-block-variables}. The speed/heading set and
the east/north component set cannot be mixed.

\begin{table}[htbp]
  \centering
  \caption{Wind data-block variables.}
  \label{tab:wind-data-block-variables}
  \begin{tabularx}{0.94\linewidth}{@{}>{\ttfamily}lX>{\ttfamily}lX@{}}
    \toprule
    \normalfont Variable & Description & \normalfont Variable & Description \\
    \midrule
    windv & Total horizontal wind speed & winde & East wind component \\
    windh & Horizontal wind heading & windn & North wind component \\
    windd & Downward wind component & windd & Downward wind component \\
    \bottomrule
  \end{tabularx}
\end{table}

Wind values may be constants or table identifiers in parentheses. TAOS uses the
parentheses to distinguish a table name from a constant, so a profile that varies with
altitude, time, or another permitted state variable can be entered through tables.

\begin{lstlisting}[style=taosmanual]
*wind geodetic windv=(winds) windh=90 windd=0
\end{lstlisting}

This interpolates horizontal speed from table \tableid{winds}, specifies wind from
the east, assigns no vertical component, and defines the components in the local
geodetic horizon coordinate system.
